% =============================================================================================================
% Skeleton of the affine-forms paper. Target: International Journal of Electrical Power & Energy Systems (IJEPES).
% Compile on Overleaf with the Elsevier "elsarticle" class (not installed in the repository container).
% Red notes (\plan, \nota) are for the authors and must disappear before submission: see paper/README.md.
% Every number in the paper must come from RESULTS.md / results/*.csv; the section and CSV are given next to it.
% =============================================================================================================
\documentclass[preprint,12pt]{elsarticle}
\usepackage{amsmath,amssymb,amsthm}
\usepackage{graphicx}
\usepackage{booktabs}
\usepackage{xcolor}
\usepackage[hidelinks]{hyperref}

\newcommand{\plan}[1]{\textcolor{red}{[PLAN: #1]}}
\newcommand{\nota}[1]{\textcolor{red}{[NOTA: #1]}}
\newcommand{\src}[1]{\textcolor{blue}{[#1]}}          % where the number comes from (RESULTS.md section / CSV)
\newtheorem{proposition}{Proposition}
\newtheorem{assumption}{Assumption}

\journal{International Journal of Electrical Power \& Energy Systems}

\begin{document}
\begin{frontmatter}

% \nota{Título: tres opciones; elegir una.}
% (a) Coverage-guaranteed probabilistic bounds on voltages and branch currents from sparse affine forms
% (b) Why linearized probabilistic power flow fails in the tails, and a second-order range that fixes it
% (c) Certified probabilistic voltage and branch-current bounds with the range of the second-order term
\title{Coverage-guaranteed probabilistic bounds on voltages and branch currents from sparse affine forms and the range of the second-order term}

\author[pucv]{Eva Ponce-Fernández}
\author[pucv]{Juan Geraldo-Ulloa}
\author[pucv]{Ignacio Araya\corref{cor}}
\cortext[cor]{Corresponding author}
\affiliation[pucv]{organization={Escuela de Ingeniería Informática, Pontificia Universidad Católica de Valparaíso},
            city={Valparaíso}, country={Chile}}
% \nota{Verificar en la guía de IJEPES el formato de autores, correo de contacto y si la revisión es simple o doble ciego.}

\begin{abstract}
% \nota{Borrador. Unas 190 palabras como en IJEPES 2026 (Basirat et al.); verificar el límite en la guía. Cifras: §28.8, §28.12.}
Linearized probabilistic power flow models are fast and widely used to bound bus voltages and branch flows under
uncertain injections, but their bounds are rarely checked where operational decisions are taken: in the extreme
tails. We show on four test systems of up to 2869 buses that Gaussian, cumulant and point-estimate methods leave
between 0.3\% and 100\% of nonlinear AC samples outside bounds designed for 0.27\%, and that correcting the
linearization with the mean and variance of the second-order term, as existing second-order methods do, is not
enough. The reason is that the linearization error is skewed and heavy-tailed. We propose to bound each output with
its sparse affine form, a distribution-free concentration radius, and the range of the exact second-order term,
sampled from a closed-form quadratic model without solving any power flow. An order-statistics argument turns the
sampled range into a provable coverage level for the second-order model. The bounds cover the AC samples in every
tested case, from 80\% load to 95\% of the loadability limit, with dispersion levels up to 30\%, correlated sources
and eight bounded input distributions, for bus voltages, branch voltage drops and branch currents, in about 10\,s on
PEGASE-2869.
\end{abstract}

\begin{highlights}
% \nota{Elsevier: normalmente 3 a 5 viñetas de hasta 85 caracteres. Verificar en la guía.}
\item Linearized probabilistic bounds fail in the tails against the AC power flow
\item The linearization error is skewed and heavy-tailed: moments cannot fix it
\item The range of the exact second-order term restores coverage
\item Order statistics give a provable coverage level at small extra width
\item Validated on four systems up to 2869 buses, for voltages and branch currents
\end{highlights}

\begin{keyword}
Probabilistic power flow \sep Affine arithmetic \sep Linearization error \sep Concentration inequalities \sep
Branch currents \sep Coverage guarantee
\end{keyword}
\end{frontmatter}

% \nota{Nomenclatura: IJEPES 2026 la incluye. Armarla al final.}

% =============================================================================================================
\section{Introduction}\label{sec:intro}

\plan{P1. Motivación. Reutilizar y acortar el primer párrafo de caja-elipse.tex (incertidumbre renovable;
necesidad de cotas de voltaje y de carga de líneas). Sin la parte de ``rigorous containment''.}

\plan{P2. Dos familias: (i) intervalos y aritmética afín (Vaccaro 2010; Dimitrovski 2004): deterministas, sufren
el problema de dependencia o se vuelven muy anchos; (ii) flujo probabilístico: Monte Carlo (caro), cumulantes
\cite{zhang2004cumulant}, estimación por puntos \cite{hong1998pem}, caos polinomial, y restricciones de azar con
linealización en el punto AC \cite{roald2018ccacopf,lubin2019chance}. Los métodos analíticos se validan comparando
momentos o curvas de densidad, casi nunca la cobertura de los cuantiles extremos.}

\plan{P3. Segundo orden: existe desde \cite{brucoli1985quadratic}; versión analítica moderna en
\cite{sandoval2022second}; no lineal por holomorphic embedding \cite{liu2018holomorphic}; expansiones para la
densidad \cite{basirat2026soeae}; estimación por puntos para las colas \cite{singh2025pem}. Todos propagan momentos
o cumulantes, validan con curvas y llegan a 118 barras como máximo (RELATED\_WORK.md, sexta pasada).}

\plan{P4. El hueco: nadie verifica (ni garantiza) la cobertura al nivel del 0,27\,\% a escala de transmisión, y
mostramos que los momentos del segundo orden no alcanzan. Concentración sobre formas afines existe fuera de
sistemas de potencia \cite{bouissou2016affine}.}

\plan{P5. Contribuciones (lista):
(1) evidencia de que los métodos linealizados (3$\sigma$, cumulantes, estimación por puntos) y la corrección de
segundo orden por momentos no cubren en las colas, y de por qué (asimetría y cola pesada del error de
linealización) [§27.1, §27.2, §28.8];
(2) cotas por forma afín dispersa + radio de concentración + rango del término de segundo orden, muestreado de un
modelo cuadrático cerrado, sin flujos de potencia [§28.1, §28.3];
(3) nivel de cobertura demostrable para el modelo de segundo orden por estadísticos de orden [§28.12];
(4) la misma construcción para caídas de voltaje y corrientes de rama, donde la corrección es imprescindible
[§28.9];
(5) validación en 4 sistemas hasta 2.869 barras, estrés de carga, correlación, distribuciones no uniformes, y
límites medidos (peor caso en vértices, cobertura simultánea) [§28.2, §28.6, §28.7, §28.10].}

\plan{P6. Organización del artículo.}

% =============================================================================================================
\section{Affine forms of the power-flow outputs}\label{sec:forms}

\subsection{Uncertainty model}
\plan{Reutilizar §II-B de caja-elipse.tex: forma afín $\hat x = x_0 + \sum_k x_k\varepsilon_k$; una fuente por $P_d$ y
otra por $Q_d$ de cada carga, con semiancho $r_k = \delta|S_k^0|$. Supuestos explícitos (Assumption~\ref{as:sym}).}

\begin{assumption}\label{as:sym}
The symbols $\varepsilon_k$ are independent, have zero mean, take values in $[-1,1]$ and have variance at most
$\sigma_\varepsilon^2$ ($\sigma_\varepsilon^2 = 1/3$ by default). Their joint law is known well enough to sample it.
\end{assumption}
\nota{Correlación: con factores compartidos $\varepsilon_k = w\zeta_{z(k)} + (1-w)\eta_k$ el supuesto se cumple sobre los
símbolos $(\zeta,\eta)$ [§28.6]. Discutirlo aquí o en §\ref{sec:limits}.}

\subsection{Sparse first-order forms}
\plan{Reutilizar §II-A y §III-A de caja-elipse.tex (linealización en el punto AC convergido, LU dispersa con COLAMD,
SuperLU, sin invertir $J$), quitando todo lo de la elipse. Escribir la forma afín de cada objetivo: barra
$\Delta|V_i| = c_i^\top\varepsilon$; caída $c_i - c_j$; corriente
$\Delta|I_\ell| = \mathrm{Re}\big(\overline{u_\ell}\,(Y_b E)_\ell\big)\varepsilon$ con $u_\ell = I_{\ell,0}/|I_{\ell,0}|$.
Subrayar que la dependencia entre barras está en los símbolos compartidos: por eso la caída sale angosta [§28.1].}

\subsection{Exact second-order term}
\plan{Coordenadas rectangulares: las ecuaciones son cuadráticas; $D = C\varepsilon - J^{-1}q(D)$;
$\Delta|V_i| = c_i^\top\varepsilon + \varepsilon^\top H_i\varepsilon + O(\delta^3)$, con $H_i$ sin formar: un vector
adjunto por barra (second\_order.py). Estado de segundo orden $\Delta V_1 + \Delta V_2$,
$\Delta V_2 = -J^{-1}q(\Delta V_1)$, para las corrientes. Citar \cite{brucoli1985quadratic,sandoval2022second}:
el modelo no es nuevo; lo nuevo es cómo se usa.}

% =============================================================================================================
\section{Why moments are not enough}\label{sec:why}
\plan{Sección corta, motivadora, con una figura. (a) A lo largo de una barra el término cuadrático es chico frente a
$\sigma$, pero en las colas desplaza la cota; (b) su distribución es asimétrica y de cola pesada (asimetría
cercana a $\pm2{,}7$, curtosis cercana a 13 en las ramas peores, diagnóstico de B4, sin CSV: generar la figura y su
CSV); (c) corrección por media y varianza: falla en 15 de 25 celdas de PEGASE [§27.2,
results/b4\_sweep.csv].}
\plan{Figura~1: histograma del término cuadrático en una barra o caída típica de PEGASE frente a la normal de igual
media y varianza; marcar el cuantil 99,865\,\%. Script por hacer: experiments/figures\_paper.py.}

% =============================================================================================================
\section{Bounds with the range of the second-order term}\label{sec:method}

\subsection{Distribution-free radius of the linear part}
\plan{Bernstein con soporte y varianza; Hoeffding con solo soporte \cite{hoeffding1963,boucheron2013concentration};
tope por la caja afín $\sum_k|c_k|$. Explicar por qué no Chebyshev ni desigualdades para unimodales: quedan topadas
por la caja afín [§28.4, results/b5\_radii.csv].}
\begin{equation}
r_\alpha = \frac{ML}{3} + \sqrt{\Big(\frac{ML}{3}\Big)^2 + 2\sigma^2 L}, \qquad L = \ln\frac{2}{\alpha},\quad
M = \max_k |c_k|,\quad \sigma^2 = \sigma_\varepsilon^2\textstyle\sum_k c_k^2 .
\end{equation}

\subsection{Sampled range of the second-order term}
\plan{Bound $[\,-r + \min(0,q_{\min}),\ r + \max(0,q_{\max})\,]$ con $q_{\min},q_{\max}$ extremos de $n$ muestras
del término cuadrático (sin flujos de potencia). Una sola pasada para todos los objetivos (NonlinearSampler).}

\subsection{A provable coverage level}
\begin{proposition}\label{prop:cert}
Under Assumption~\ref{as:sym}, let $Y = L + Q$ be a target of the second-order model, with $L = c^\top\varepsilon$
and $Q$ its quadratic part. Let $q_{\min}, q_{\max}$ be the extremes of $n$ independent samples of $Q$, drawn
independently of $\varepsilon$, and let $r$ satisfy $P(|L| > r)\le\alpha_L$. Then
\[
P\Big(Y \notin \big[-r + \min(0,q_{\min}),\; r + \max(0,q_{\max})\big]\Big) \;\le\; \alpha_L + \frac{2}{n+1}.
\]
\end{proposition}
\begin{proof}
If $|L|\le r$, $Q\le q_{\max}$ and $Q\ge q_{\min}$, then $Y\le r + \max(0,q_{\max})$ and
$Y\ge -r + \min(0,q_{\min})$. Hence the event in the statement is contained in
$\{|L|>r\}\cup\{Q>q_{\max}\}\cup\{Q<q_{\min}\}$. The new $Q$ and the $n$ samples are exchangeable, so the probability
that the new one is the strict maximum of the $n+1$ values is at most $1/(n+1)$; the same holds for the minimum.
The union bound gives the result.
\end{proof}
\plan{Corolario: con $\alpha_L = \alpha - 2/(n+1)$ la cota tiene nivel $\alpha$; con $n = 4000$ el ancho sube 2{,}2\,\%
en mediana [§28.12, results/b11\_certified.csv]. Aclarar que la garantía es para el modelo de segundo orden; el
tercer orden se mide contra AC.}

\subsection{Algorithm and cost}
\plan{Pseudocódigo de 6 pasos: (1) punto AC; (2) LU y formas afines; (3) modelo de segundo orden; (4) $n$ muestras
de los términos no lineales en una pasada; (5) radios; (6) cotas. Tabla de costos [§28.12,
results/b11\_timing.csv]: 6,7\,s voltajes, 10,2\,s con corrientes en PEGASE-2869.}

% =============================================================================================================
\section{Case studies}\label{sec:cases}

\subsection{Setup}
\plan{Sistemas: IEEE-33, IEEE-118, PEGASE-1354, PEGASE-2869 (pandapower y PGLib \cite{pglib2019,thurner2018pandapower}).
Fuentes: todas las cargas, $P$ y $Q$ separadas. Celdas: 80\,\% a 95\,\% del límite de cargabilidad,
$\delta$ de 5 a 30\,\%. Validación: 20.000 muestras AC por celda (solver por lotes, tolerancia $10^{-9}$),
métricas: peor objetivo, objetivos sobre 0,27\,\%, fracción de muestras con algún objetivo fuera, ancho.
Software y versiones (README.md).}

\subsection{Linearized bounds fail in the tails}
\plan{Tabla: 3$\sigma$, cumulantes + Cornish-Fisher, estimación por puntos + Cornish-Fisher, Bernstein sin corregir
y con el rango, en las 3 celdas de los 4 sistemas [§28.8, results/b8\_ppf\_baselines.csv].}

\subsection{Coverage across loading and dispersion}
\plan{Figura~2: mapa carga × $\delta$ de la peor violación para 3$\sigma$, momentos y rango (PEGASE-2869)
[§27.2, §28.3, results/b4\_sweep.csv, b5\_sweep.csv]. Tabla con los dos sistemas nuevos [§28.7, b7\_generality.csv].}

\subsection{Branch voltage drops and branch currents}
\plan{Tabla: caídas (forma afín frente a restar intervalos, 4 a 21 veces) [§28.1, b5\_mc.csv, b5\_width.csv];
corrientes (error lineal hasta 10,7\,p.u.; sin el rango, hasta 90\,\% de violación; con el rango, 0) [§28.9,
b9\_branch\_currents.csv]. Figura~3: ancho de la cota de corriente frente a $|I_0|$.}

\subsection{Correlated sources and non-uniform distributions}
\plan{[§28.6, b6\_correlated.csv]: factores compartidos cubren, la concentración se pierde con un factor común
fuerte, la cópula mal especificada falla, la slack sola falla con correlación fuerte. [§28.10,
b10\_marginals.csv]: con media cero y soporte acotado cubren; fuera del soporte, nada.}

\subsection{Certified level and computational cost}
\plan{[§28.12, b11\_certified.csv, b11\_timing.csv]. Comparar con Monte Carlo AC por lotes (60\,s) y con pandapower
en lazo (unos 3.800\,s) [C8].}

% =============================================================================================================
\section{Limits}\label{sec:limits}
\plan{Lista honesta, con su respaldo:
(a) no es una garantía determinista: el peor caso de segundo orden no encierra la solución AC en los vértices a
escala de transmisión [§28.2]; (b) la cobertura simultánea sin supuesto de distribución cuesta lo mismo que el peor
caso [§28.7]; (c) la ventaja sobre el peor caso depende de cuántas fuentes mueven cada salida (IEEE-118: ninguna)
[§28.7]; (d) la correlación debe entrar como estructura de factores; con la slack sola y correlación fuerte falla
la linealización [§28.6]; (e) el soporte acotado es esencial [§28.10]; (f) tercer orden medido, no acotado; (g)
sin límites de reactivos de los generadores; (h) sin datos reales de carga.}

\section{Conclusions}\label{sec:conclusion}
\plan{Cuatro frases: qué falla, por qué, qué lo arregla, con qué garantía y a qué costo.}

% =============================================================================================================
\section*{CRediT authorship contribution statement}
\nota{Completar (IJEPES lo exige; ver el formato en Basirat et al. 2026).}

\section*{Declaration of competing interest}
The authors declare that they have no known competing financial interests or personal relationships that could
have appeared to influence the work reported in this paper.

\section*{Declaration of generative AI and AI-assisted technologies}
\nota{Elsevier pide declarar el uso de IA generativa en la redacción. Este trabajo usó un asistente de IA para el
código, los experimentos y borradores; el grupo debe decidir y redactar la declaración según la guía vigente.}

\section*{Data availability}
\nota{Propuesta: código y resultados en un repositorio público con DOI (por ejemplo Zenodo); los casos de prueba
son públicos (pandapower, PGLib, MATPOWER).}

\bibliographystyle{elsarticle-num}
\bibliography{refs}
\end{document}
